\documentclass{vskpou}

\usepackage[utf8]{inputenc}
\usepackage[IL2]{fontenc}
\usepackage{graphicx}
\usepackage{tabularx} % pro lepší zalamování v tabulkách

\jazyky{czech}{english}
\hlavnijazyk{czech}

\autor{Jakub Papánek}
\vedouci{RNDr. Marek Vajgl, Ph.D.}
\rok{2025}
\fakulta{Přírodovědecká Fakulta}{Faculty of Science}
\katedra{Katedra Informatiky}{Department of Computer Science}
\nazev{Tvorba aplikace pro pokročilé testování}{Thesis Title}
\typprace{Bakalářská Práce}{Bachelor Thesis}

\usepackage{todonotes}
\newcommand{\mv}[1]{\todo[inline,color=red!50]{MV: #1}}
\newcommand{\jp}[1]{\todo[inline,color=blue!25]{PK: #1}}

\pohlavi M
\prohlaseni{V Ostravě dne 16. 12. 2025}

\podekovani{...}

\abstrakt{
Cílem této práce je navrhnout a implementovat full-stack webovou aplikaci pro správu testů, inspirovanou systémem Moodle, avšak s důrazem na moderní uživatelské rozhraní a specifické analytické funkce pro sebetestování.
}{
The goal of this thesis is to design and implement a full-stack web application for test management, inspired by the Moodle system, but with an emphasis on a modern user interface and specific analytical functions for self-testing.
}

\klicovaslova{testování, webová aplikace, React, Spring Boot, vzdělávání}{testing, web application, React, Spring Boot, education}
\resume{Resumé v češtině...}{Summary in English...}

\begin{document}

\tableofcontents 

\section{Osnova práce}

\mv{Zdravím, do téhle sekce vám napíši, jak bych si představoval osnovu práce}
\jp{Pomocí "\textbackslash jp\{...\}" dělejte svoje vlastní komentáře}

Základní kapitoly:

\begin{itemize}
    \item Úvod -- představení toho, co se bude v práci řešit. Je to info pro čtenáře, kterým ho chcete nalákat na to, co bude číst. Takže rychlé představení současného stavu a naznačeného způsobu řešení
    \item Cíle práce -- představení cílů práce (čeho se má dosáhnout) a základního způsobu, který použijete k dosažení cílů práce
    \item Analýza současného stavu -- představení jednak metodiky, pomocí které budete hledat aktuální řešení (tj. jak jsme se bavili: jaké jsou kritéria, jakým způsobem se budou hodnotit), druhak představení současně dostupných nástrojů pro dané řešení (stručné předsdtavení + odkaz na nástroj, ohdonocení kritérií podle metodiky) + na závěr shrnutí, zda lze nějaký existující nástroj použít a není třeba tvořit nic nového, nebo zda se tedy bude dělat vlastní řešení
    \item Metodika řešení práce - představení postupů, přístupů a nástrojů, které použijete při řešení práce; jakým způsobem se bude postupovat při sběru požadavků, jak se budou proritizovat, že/zda bude vize, risk-list, git, jak testování atp.
    \item Odteď dále kapitoly ohledně vlastního řešení práce. Je mi jedno, jak to rozčleníte -- nejdřív architektura a pak po kapitolách BE/FE/DB, nebo po funkcionalitách, ale nějak smysluplně to pojmout
    \item Představení vytvořeného řešení -- rychlé představení vytvořeného řešení a toho, jakt o vypadá. je to o to "prodat co vzniklo", ne o tvorbě uživatelské příručky
    \item Shrnutí -- zde sepíšete, zda/jak se podařilo splnit cíle práce, co se povedlo co se nepovedlo + případně co by šlo dělat do budoucna
    \item Závěr -- zde napište, co vlastně vzniklo a jak to dopadlo.
\end{itemize}

V průběhu textu nezapomeňte přidávat reference (pořád si je průběžně ukldádejte, ať je pak nemusíte hledat).

Připravte si na tohle všechno vlastní kapitoly, případný existující text do toho přesuňte a dejte echo, že je hotovo a můžu kontrolovat, díky.

\section{úvod}

\subsection{Současný stav problematiky}

\subsection{řešení problému}


\section{cíle práce}

\subsection{cíle práce}

\subsection{způsob dosažení cíle}
\jp{způsob ? programování ? nevím co zde patří}


\section{analýza součaného stavu}
\jp{ještě se bude upravovat}

\subsection{metodika analýzy existujících řešení}

\subsubsection{základní funkcionality}

\begin{table}[h]
\centering
\begin{tabular}{|l|l|}
\hline
\textbf{základní funkcionality} & \textbf{váha} \\ \hline
dostupnost & 9 \\ \hline
vytváření testů & 5 \\ \hline
otázky & 5 \\ \hline
absolvování testů & 8 \\ \hline
analytika a statistiky & 7 \\ \hline
UX/UI & 3 \\ \hline
\end{tabular}
\end{table}

\subsubsection{hodnocení}

\begin{itemize}
    \item jednotlivé základní funkcionality jsou hodnoceny A-E + váha
    \item hodnocení se potom budou pomocí vzorce agregovat do finalního hodnocení
    \item každá základní funkcionalita bude mít několik definovaných konkrétních kritériích, každé kritérium má určitý počet bodů který může dosáhnout podle toho jestli je splněno nebo ne
    \item body se sečtou a podle toho se udělí známka celé základní funkcionality
    \begin{itemize}
        \item 0-1 bod $\Rightarrow$ E
        \item 2 body $\Rightarrow$ D
        \item 3 body $\Rightarrow$ C
        \item 4 body $\Rightarrow$ B
        \item 5 bodů $\Rightarrow$ A
    \end{itemize}
\end{itemize}

\subsubsection{Dostupnost (Accessibility \& Hosting)}

Hodnotíme, jak snadno se k nástroji dostane koncový uživatel (student) bez pomoci instituce.

\begin{table}[h]
\centering
\begin{tabular}{|p{0.7\textwidth}|l|}
\hline
\textbf{Kritéria} & \textbf{Známka} \\ \hline
Veřejně dostupná SaaS. Uživatel si založí účet (zdarma nebo s jasným ceníkem) a může okamžitě využívat všechny klíčové funkce bez nutnosti cokoliv instalovat nebo hostovat. & A (5 bodů) \\ \hline
Veřejně dostupná SaaS, ale s výraznými paywally (např. zamknuté základní typy otázek za prémiové předplatné) nebo s velkým omezením počtu testů ve free verzi. & B (4 body) \\ \hline
Self-hosted open-source (např. Moodle). Je zcela zdarma a otevřený, ale vyžaduje, aby si uživatel sám zařídil server, databázi a provedl instalaci (velká technická bariéra). & C (3 body) \\ \hline
Vyžaduje příslušnost k organizaci (škole), ale po získání přístupu umožňuje studentovi plně samostatnou práci a tvorbu vlastních materiálů. & D (2 body) \\ \hline
Uzavřený B2B/Enterprise systém. Nelze se do něj registrovat jako jednotlivec a bez učitele/školy je naprosto nedostupný a nepoužitelný. & E (1 bod) \\ \hline
\end{tabular}
\end{table}

\begin{itemize}
    \item zde hodnocení funguje trochu jinak protože každé kritéria mají jinou relevanci a také vždy platí jen jedno z nich (platforma nemůže být zároveň self-hosted a SaaS)
\end{itemize}

\subsubsection{Vytváření testů a organizace}

\begin{table}[h]
\centering
\begin{tabular}{|p{0.7\textwidth}|l|}
\hline
\textbf{Kritéria} & \textbf{hodnocení v bodech} \\ \hline
má/nemá časový limit & 1 bod \\ \hline
má/nemá kategorie & 1 bod \\ \hline
umožnuje vytvořit testy s jakým koliv typem otázky * & 1 bod \\ \hline
UX: celková flow, jak dobře se to používá ** & 2 body \\ \hline
\end{tabular}
\end{table}

\begin{itemize}
    \item * - to znamená že když platforma podporuje například otázky s jednou odpovědí, otázky s více odpověďmi a flashcards tak umožňuje uživatelovi vytvořit test se všemi typy otázek a ne například jen s flashcards.
    \item ** - zde chceme hodnotit hodnotit UX konkrétně protože tato část aplikace může výrazně trpět na špatný UX
    \item ** - UX má 2 body protože je to víc abstraktní kritérium proto je vyjádření na stupnici s rozpětím víc efektivní
    \item vynecháváme základní kritéra jako title, description… atd. protože to jsou základní funkce u kterých se očekává že jsou pokryté na každé platformě proto je nemá smysl zakomponovat do hodnocení
\end{itemize}

\subsubsection{Typologie a obsah otázek}

\begin{table}[h]
\centering
\begin{tabular}{|p{0.7\textwidth}|l|}
\hline
\textbf{Kritéria} & \textbf{hodnocení v bodech} \\ \hline
má/nemá otázku typu jedna možnost & 0,5 bodu \\ \hline
má/nemá otázku typu více možností & 0,5 bodů \\ \hline
má/nemá otázku typu textová odpověď & 0,5 bodů \\ \hline
má/nemá otázku typu pravda/nepravda & 0,5 bodů \\ \hline
má/nemá otázku typu dvojice & 0,5 bodů \\ \hline
má/nemá otázku typu Seřazení & 0,5 bodů \\ \hline
má/nemá otázku typu Flashcard * & 1 bod \\ \hline
podporuje přidání obrázků k otázkám & 1 bod \\ \hline
\end{tabular}
\end{table}

\begin{itemize}
    \item pro platformu je klíčové aby podporovala co nejvíc typů otázek aby se odlišila od ostatních platforem proto chceme hodnotit všechny
    \item * - Flashcard má větší váhu protože to nepodporují všechny platformy
\end{itemize}

\subsubsection{Absolvování testů}

\begin{table}[h]
\centering
\begin{tabular}{|p{0.7\textwidth}|l|}
\hline
\textbf{Kritéria} & \textbf{hodnocení v bodech} \\ \hline
UX: možnost zobrazení výsledků hned po dokončení testu & 1 bod \\ \hline
UX: možnost opětovného spuštění testu hned po dokončení & 1 bod \\ \hline
UX: test se dá jednoduše spustit pro sebe bez toho aby se musel spustit přes nasdílený odkaz v druhém okně * & 0 nebo 2 body \\ \hline
umožňuje volnou navigaci mezi otázkami & 1 bod \\ \hline
\end{tabular}
\end{table}

\begin{itemize}
    \item * - hodnocení nemůže dosáhnout 1 bodu, buď 0 bodů nebo 2, hodnocení nelze vyjádřit na stupnici
    \item zde chceme hodnotit UX konkrétně protože to jsou důležité kritéria které chceme pokrýt
\end{itemize}

\subsubsection{Analytika a statistiky}

\begin{table}[h]
\centering
\begin{tabular}{|p{0.7\textwidth}|l|}
\hline
\textbf{Kritéria} & \textbf{hodnocení v bodech} \\ \hline
má/nemá historii pokusů & 1 bod \\ \hline
má/nemá jakékoliv statistické vyjádření úspěšnosti testu * & 0 nebo 2 body \\ \hline
má/nemá dlouhodobé vedení statistik ** & 1 bod \\ \hline
instantní zpětná vazba ukonceni testu (revize, vyjádření úspěšnosti) & 1 bod \\ \hline
\end{tabular}
\end{table}

\begin{itemize}
    \item * - hodnocení nemůže dosáhnout 1 bodu, buď 0 bodů nebo 2, hodnocení nelze vyjádřit na stupnici
    \item ** - není to tak že jsou statistiky dostupné jen po dokončení jednoho testu ale jsou dlouhodobě vedené skrze všechny pokusy
\end{itemize}

\subsubsection{UX/UI}

\begin{table}[h]
\centering
\begin{tabular}{|p{0.7\textwidth}|l|}
\hline
\textbf{Kritéria} & \textbf{hodnocení v bodech} \\ \hline
celkové UI & 2,5 bodů \\ \hline
celkové UX & 2,5 bodů \\ \hline
\end{tabular}
\end{table}

\begin{itemize}
    \item tyto kritéria mají hodnocení až 2,5 bodů protože jsou víc abstraktní proto je vyjádření na stupnici s rozpětím víc efektivní
\end{itemize}

\subsubsection{agregace finálního hodnocení}

\begin{enumerate}
    \item hodnocení základních funkcionalit platformy (A-E) se převede na body A=5, B=4, C=3, D=2, E=1 $\Rightarrow$ \textit{Hi}
    \item Každé základní funkcionalitě přiřadíme váhu \textit{Wi} (viz. tabulka nahoře)
    \item Finální známka pro danou platformu se vypočítá pomocí tohoto vzorce:\\
    $S = (w_1 \cdot H_1 + w_2 \cdot H_2 + \dots + w_n \cdot H_n) / (w_1 + w_2 + \dots + w_n)$
\end{enumerate}

\begin{figure}[htbp]
    \centering
    \includegraphics[width=0.8\textwidth]{obrazek.png}
    \caption{Agregace finálního hodnocení}
\end{figure}
\jp{obrázek dostane bílé pozadí}

\subsection{představení a analýza existujících řešení}

\subsection{shrnutí analýzy existujících řešení}


\section{metodika řešení práce}
\subsection{nástroje}
\jp{u každého nástroje bude jen v pár větách popsané co dělají a jak jsou užitečné}
\subsubsection{vscode}

\subsubsection{vscode rozšíření: \textbf{Database Client}}
\jp{toto mi příjde relevantní protože jsem to používal jako GUI pro zobrazení obsahu databáze}
\subsubsection{git}

\subsubsection{jira}

\subsubsection{exceli draw}

\subsubsection{draw.io}

\subsubsection{overleaf}

\subsubsection{notion (?)}
\jp{k shromažďování informací a k dokumentaci, měl bych to tam dát?}
\subsection{postup vývoje}
\jp{napíšu o tom že to byl scrum s 1 týdenními sprinty}
\subsection{sběr požadavků}
\jp{co se myslí "jakým způsobem se bude postupovat při sběru požadavků" ? a "jak se budou proritizovat" ?}

Požadavky budou sbírány tak aby aplikace primárně podporovala základní funkcionality k tomu aby mohly být implementovány hlavní funkcionality které jsou cílem této práce.
\jp{potom se tam ještě dopíše že bude vize a risk-list}

\section{řešení práce}

Vývoj celé aplikace byl realizován s využitím pestré palety moderních softwarových nástrojů a knihoven, které zajišťují plynulý chod, bezpečnost a snadnou udržovatelnost. Projekt je striktně rozdělen na klientskou a serverovou část.

\subsection{frontend}

Aplikace je navržena jako Single Page Application (SPA), což je pro takto vysoce interaktivní systém naprosto nezbytné. Architektura SPA zajišťuje plynulý chod bez znovunačítání celé stránky, což je kritické pro funkce jako odpočet časovače, plynulé přeskakování mezi otázkami, Drag \& Drop. React představuje pro tento účel ideální volbu díky své schopnosti efektivně spravovat složitý stav klientského rozhraní. Zásadním benefitem je obrovský ekosystém komunitních balíčků (packages), který je ze všech javascript frameworků/knihoven nejrozsáhlejší.

\subsection{backend}

Pro backend je vynikající volbou framework Spring Boot, jenž poskytuje základ pro tvorbu RESTful API, správu komplexní business logiky pro hodnocení pokusů a bezpečně řízenou bezstavovou autentizaci. Architektura aplikace z tohoto frameworku výrazně těží, a to díky modulu Spring AI, který umožňuje přímou a efektivní integraci s velkými jazykovými modely (LLM). 

\subsection{Architektura systému}

Celá fullstack aplikace je postavena React, Springboot a MySQL. Aplikace je navržená na principech vrstvené architektury (tzv. n-layer), která odděluje jednotlivé části systému do specifických vrstev, jako jsou Frontend + REST kontrolery, vrstva business logiky (service) a datové entity. Tato volba je praktická, protože projekt je vyvíjen jedním člověkem. Vrstvená architektura je díky své a povědomosti velmi jednoduchá na pochopení i implementaci. Vzhledem k tomu, že se u této aplikace neočekává dlouhodobé či časté nasazování do produkčního prostředí a zároveň se nepočítá s nasazením pro masivní počet uživatelů, nejsou nedostatky vrstvené architektury problémem. Typické nevýhody, jako je například obtížnější horizontální škálovatelnost nebo potenciální výkonnostní úzká hrdla při obrovské zátěži, se v tomto konkrétním rozsahu vůbec neprojeví.

\begin{itemize}
    \item \textbf{Prezentační vrstva (Frontend + API Endpointy)}
    \begin{itemize}
        \item Skládá se z klientského frontendu a API endpointů na backendu. Uživatelské rozhraní frontendu je tvořeno react komponentami. Uživatelské rozhraní je zcela responzivní. Komunikace s backendem probíhá přes REST API endpointy, u kterých je vyžadována platná autentizace pomocí JWT tokenu (s výjimkou několika málo veřejných cest pro registraci a přihlášení).
    \end{itemize}
    
    \item \textbf{Aplikační vrstva (větčina backend aplikační logiky)}
    \begin{itemize}
        \item Tvoří funkční jádro celého systému a je v ní soustředěna většina hlavní aplikační logiky. Tato vrstva zpracovává požadavky přicházející z prezentační vrstvy, je doménově dělena a řídí klíčové procesy, jako je správa uživatelských účtů, vytváření kvízů, organizace kategorií nebo zpracování odevzdaných pokusů.
    \end{itemize}
    
    \item \textbf{Servisní vrstva}
    \begin{itemize}
        \item Funguje jako specializovaná pomocná vrstva, kterou aplikační vrstva využívá k delegování složitějších operací. Obsahuje nezbytné pomocné servisy a komponenty určené pro specifické úkoly. Patří sem služba pro vyhodnocení kvízu a také dedikovaná služba pro dotazování jazykového modelu (LLM), jež se stará o automatické posuzování volných textových odpovědí.
    \end{itemize}
    
    \item \textbf{Databázová vrstva (konektor, ORM, Databáze)}
    \begin{itemize}
        \item zajišťuje uložení všech dat do relační databáze MySQL. Datová integrace s backendem je řešena za využití frameworku Hibernate, který slouží jako nástroj pro objektově-relační mapování (ORM), a technologie Spring Data JPA, jež usnadňuje manipulaci s databázovými entitami. Pro samotné spojení aplikace s databází slouží JDBC konektor.
    \end{itemize}
\end{itemize}

\subsection{přehled použitých technologií}

\begin{table}[h]
\centering
\begin{tabular}{|p{0.35\textwidth}|p{0.6\textwidth}|}
\hline
\textbf{Logický celek} & \textbf{Použité technologie a knihovny} \\ \hline
\textbf{Jazyky a hlavní frameworky} & TypeScript, React, Java, Spring Boot \\ \hline
\textbf{Build nástroje a kontejnerizace} & Vite, Maven, Docker, Docker Compose \\ \hline
\textbf{Autentizace a bezpečnost} & Spring Security, JWT \\ \hline
\textbf{Databáze a ORM} & MySQL, Spring Data JPA, Hibernate \\ \hline
\textbf{Umělá inteligence} & Spring AI, Ollama, llama3.1 \\ \hline
\textbf{Uživatelské rozhraní a styly} & Tailwind CSS, shadcn/ui, Radix UI \\ \hline
\textbf{Správa stavu, routing a HTTP} & Redux Toolkit, React Redux, React Router DOM, Axios \\ \hline
\textbf{Grafy a pokročilé interakce} & ECharts, Recharts, dnd-kit \\ \hline
\end{tabular}
\end{table}


\section{Představení řešení}

\section{shrnutí}

\section{závěr}

\end{document}